\chapter*{About This Book}
\addcontentsline{toc}{chapter}{About This Book}
\markboth{About This Book}{About This Book}

\emph{Large Language Models in Finance} is written for a mixed audience of academic
researchers, graduate students, and industry practitioners---quantitative analysts, data
scientists, and engineers---who want both the conceptual foundations and the practical
machinery to deploy large language models on real financial problems. It assumes comfort
with probability, linear algebra, and Python, but it does not assume prior exposure to
modern deep learning: the necessary background is built up from first principles.

\section*{What you will find}

The project has three companion components that share a common chapter numbering:

\begin{itemize}
  \item \textbf{The book} (this volume) --- the deepest treatment, with definitions,
        derivations, and worked financial examples.
  \item \textbf{Course materials} --- lecture notes, slides, and exercises that track a
        subset of the book for classroom use.
  \item \textbf{Companion code} --- under \texttt{code/notebooks/<chapter>/}: a
        \texttt{demo.ipynb} that reproduces the chapter's worked figure, an
        \texttt{exercises.ipynb} lab, and deterministic \texttt{gen\_*.py} figure
        generators. Every figure in the book can be regenerated by running
        \texttt{code/run\_illustrations.sh}; the generators are network-free and seeded so
        the figures reproduce exactly.
\end{itemize}

\section*{How to read it}

Each chapter is layered so that two kinds of reader can use it. Material in a
\textsf{context} box gives the big picture---the intuition, the financial motivation, and
the practical takeaway---and can be read on its own. Material in a \textsf{deepdive} box
goes under the hood---the formal definitions, proofs, and implementation details---and can
be skipped on a first pass without losing the thread. Definitions, examples, remarks, and
illustrations are called out as labelled environments throughout.

End-of-chapter exercises are tagged by difficulty: \textbf{[B]} beginner, \textbf{[I]}
intermediate, and \textbf{[A]} advanced. Cross-references use clickable links, so a pointer
such as \Cref{ch:intro} can be followed directly in the PDF.

\section*{How the chapters build}

The early chapters establish foundations---what financial text is and how it is
represented, the architecture and training of large language models, and how to query them
in practice. The middle chapters apply those foundations to specific finance tasks:
domain-specific models, agents, sentiment, summarization and information extraction,
valuation, credit risk, portfolio construction and trading, and regulatory technology. The
final chapters address the cross-cutting concerns of any production deployment:
explainability, the evaluation of model limitations, broader applications and frontier
research, and privacy-preserving local deployment. The appendices are hands-on guides to
the tooling---Claude Code, Codex CLI, local HuggingFace inference, commercial SDKs,
Git/GitHub, and Cline.

\section*{A note on currency}

The field moves quickly. Specific model identifiers, prices, and benchmark numbers in the
text are illustrative as of writing and will change; the \emph{techniques}---tokenisation,
fine-tuning, retrieval-augmented generation, agentic orchestration, and the governance
practices around them---are architecture-level patterns that generalise across model
versions. Where a quantitative result comes from a working paper that the reader cannot
independently verify, the text states it qualitatively rather than as an exact figure.
